\usepackage{amsmath,amssymb,hyperref}
\usepackage{an-optional-package}
\newcommand{\R}{\mathbb{R}}
\newcommand{\vect}[1]{\mathbf{#1}}
\newcommand{\inner}[2]{\langle #1,#2\rangle}
\newcommand{\important}[1]{\textbf{#1}}
\newcommand{\N}{\mathbb{N}}
\renewcommand{\N}{\mathbb{N}_0}
\DeclareMathOperator{\rank}{rank}

\section{Draft preview playground}
This is a lightweight note: it needs no document class, document environment,
or PDF build. Open it with \texttt{:Dp} and edit it to see the preview refresh.
The package declarations above are compatibility hints; no packages are executed.
Even the invented package is ignored without a warning.

\subsection{Original example}
this is some text

this are some equations
\begin{equation}
  \int_{-\infty}^{\infty} e^{-x^2}\,dx = \sqrt{\pi}.
\end{equation}

\begin{itemize}
  \item First item
  \item Second item
\end{itemize}


\section{Text and links}
An ordinary paragraph can contain \textbf{bold text}, \textit{italic text},
\emph{emphasis}, \texttt{monospace code}, and \text{a plain text wrapper}.
The custom prose command makes this \important{important}.

Here is \href{https://www.latex-project.org}{a named link}, and here is a URL:
\url{https://www.latex-project.org}.

This sentence uses an explicit line break.\\The next sentence starts below it.

An unknown wrapper can still preserve its content:
\unknownwrapper{This text remains visible, including \textbf{nested formatting}.}

\section{Lists}
\begin{itemize}
  \item A simple bullet with \textbf{bold text}.
  \item A bullet with inline mathematics: $x\in\R$.
  \item A bullet containing a nested numbered list:
  \begin{enumerate}
    \item Read the note.
    \item Change an equation and watch the live preview.
    \item Use Ctrl-G to synchronize source and preview.
  \end{enumerate}
\end{itemize}

\begin{enumerate}
  \item Numbered lists work on their own, too.
  \item They can contain a nested bullet list:
  \begin{itemize}
    \item Keep the source as ordinary LaTeX.
    \item Preview a fragment without compiling a document.
  \end{itemize}
\end{enumerate}

\section{Quotation and centering}
\begin{quote}
A quotation can contain \textit{formatted prose} and inline math such as
$\inner{\vect{x}}{\vect{y}}=0$.
\end{quote}

\begin{center}
This passage is centered, with \textbf{a little emphasis}.
\end{center}

\section{Mathematics and custom commands}
Inline math works with dollar signs, $e^{i\pi}+1=0$, or parentheses,
\(a^2+b^2=c^2\).

The zero-argument commands give $\R$ and $\N$.
The one-argument command gives $\vect{x}$, and the two-argument command gives
$\inner{\vect{x}}{\vect{y}}$.
The definition of N above is replaced using renewcommand, so $0\in\N$.
The declared math operator gives $\rank A=2$.

\[
  \frac{d}{dx}\sin x = \cos x,
  \qquad \sum_{k=1}^{n} k = \frac{n(n+1)}{2}.
\]

\begin{equation*}
  A = \begin{pmatrix}1 & 0\\0 & 1\end{pmatrix},
  \qquad \vect{x}\in\R^2,
  \qquad \rank A = 2.
\end{equation*}

\[
  \begin{aligned}
    f(x) &= x^2+2x+1,\\
         &= (x+1)^2.
  \end{aligned}
\]

\section{Reading and diagnostics}
Try Ctrl-G from the source, or select some preview text and use Ctrl-G there.
Use j and k to scroll, Ctrl-D and Ctrl-U for half-page motion, and gg or G to
jump to the beginning or end. The sun/moon control changes the appearance.

To try diagnostics, add an unknown command without an argument, or an invalid
math command inside an equation. Click the small problem indicator to copy all
diagnostics, or click an individual diagnostic to copy just that message and
its source line. The clicked control briefly dims after a successful copy.

\subsection{Deliberate limits}
This preview supports a small LaTeX subset. Macro arguments must be braced;
optional defaults, TeX scoping, custom environments, and package execution are
not supported. Put list, quote, and center environments at the start of a line.
Use inline math inside lists, and place display equations outside them.
